% Summary of graf2022: Typological implications of tier-based strictly local movement

\begin{abstract}
\begin{enumerate}
  \item Analyzing movement as TSL-2 provides a unified analysis of a bunch of phenomena.
\end{enumerate}
\begin{quote}
A cognitive system that can handle basic TSL-2 movement must be able to also handle these phenomena.

\end{quote}
\begin{enumerate}[start=2]
  \item It also rules out a bunch of unattested phenomena.
  \item However, a lot of unattested phenomena are still predicted.
\end{enumerate}
\end{abstract}

\section*{Attested, predicted}

Varying the number of dependents in a Daughter-String Language:

\begin{itemize}
  \item Multiple Wh-Movement
  \begin{itemize}
    \item allow ``at least one'' f-
  \end{itemize}
  \item Expletive Constructions/Optional Movement
  \begin{itemize}
    \item allow ``at most one'' f-
  \end{itemize}
  \item The That-Trace Effect
  \begin{itemize}
    \item create a separate tier where \textit{that} can't have [nom-] daughters, and where adjuncts intervene
  \end{itemize}
  \item The Anti-That-Trace Effect
  \begin{itemize}
    \item same as above, except that a silent complementizer REL can't have [nom-] daughters.
  \end{itemize}
\end{itemize}

Using interveners on the tier:

\begin{itemize}
  \item Islands
  \begin{itemize}
    \item project whatever intervenes
  \end{itemize}
  \item Wh-Agreement in Irish
  \begin{itemize}
    \item posit two different Cs which both project onto the wh-tier. One is allowed between wh-elements, and the other isn't.
  \end{itemize}
\end{itemize}

\section*{Unattested, Unpredicted}

Weird island types:

\begin{itemize}
  \item Counting
  \begin{itemize}
    \item Gang-up islands, Cowardly islands, Rationed islands
  \end{itemize}
  \item Able to see non-locally
  \begin{itemize}
    \item Configurational islands, Cowardly islands, Narcissistic islands, Rationed islands
  \end{itemize}
  \item What a mover dominates determines its potential islands
  \begin{itemize}
    \item Discerning islands
  \end{itemize}
\end{itemize}

\section*{Unattested, Still predicted}

Languages have remarkably similar constraints.

\begin{quote}
What TSL needs is a theory of tier projections.

\end{quote}

Also, some things seem probably not TSL, but maybe could be captured with other subregular classes:

\begin{itemize}
  \item Freezing Effects
  \item the Coordinate Structure Constraint
\end{itemize}

\subsection*{great long quote about this}

\begin{quote}
languages are much more principled and systematic than the free variation account [which says that anything TSL-2 is bound to show up places] predicts.

If tier projections vary freely across tiers, languages, and speakers, why then do we find no languages that completely lack the adjunct island constraint? Why do even those languages where relative clauses do not induce island effects still show processing effects that suggest that they are islands (Tutunjian et al. 2017)? Why isn't there a language where the facts for Truswell sentences are exactly the other way around, with infinitival T opaque to extraction whereas finite T allows for it? And why isn't there an analogue of the \textit{that}-trace effects that targets objects instead of subjects? ...

This shows that TSL-2 in its current form still overgenerates and is too lax a restriction on the typology of island constraints. However, the TSL tier projection also provides a natural locus for addressing this overgeneration. \textbf{What TSL-2 needs is a theory of tier projections.} This could come in the form of substantive universals, perhaps coupled with abstract notions like monotonicity (Graf 2019, 2020, Moradi 2019, 2020, 2021). Alternatively, there may be restrictions on the relation of tiers to each other, akin to the constraints on harmony tiers identified by Aksënova and Deshmukh 2018. ...

.... As far as I can tell, some conditions on movement simply are beyond the purview of TSL-2, e.g. freezing effects and the Coordinate Structure Constraint. .... [t]here are alternative models of subregular dependencies in syntax, foremost constraints on string representations obtained from dependency trees (Graf and Shafiei 2019, Shafiei and Graf 2020) and the class of constraints recognizable by sensing tree automata (Graf and De Santo 2019).

[basically, other ways of doing subregular constraints over trees could also shore up the other spots]

\end{quote}
